\subsection{Additional controls and sensitivity analyses}
\label{app:revision-controls}

The supplementary controls used the same initialization and baseline as
the main comparisons, with grid size $L=64$, dealer depth $h=100$, and 3,000 post-intervention
steps unless otherwise specified. Every switched run was compared
with a no-switch control. Held conditions calculated full
financial histories as well as decisions. A Sell pulse wrote Sell
at the selected interior traders immediately before the first post-baseline orders,
while retaining the baseline rule for later updates. Traders already in
Sell were unchanged. A persistent intervention instead changed the
selected rule without overwriting decisions, so its first price remained
identical to the no-switch price. These paired interventions compared
the same fractions and starting markets under both technical-trader
conditions (Supplementary Table~\ref{tab:revision-pulse}). The rule
screen tested all eighteen ways of changing one response in the
baseline lookup table, using twenty starting markets per condition
and three fractions. Each screen entry was descriptive rather
than a separate confirmatory hypothesis test
(Supplementary Table~\ref{tab:revision-rules}).

The band control froze the recent-departure scale at its value immediately
before the rule switch, while each technical trader retained its multiplier
and the reference price continued to update. Adaptive and frozen runs
shared their complete pre-switch histories. The frozen widths could
therefore differ across dealer depths and starting markets. At fraction $f=0.175$,
the estimated majority-probability difference between dealer depths $h=0.1$ and
$h=100$ was -0.06 with adaptive widths and -0.26 with
frozen widths. Subtracting those two differences gave
0.20, with a 95\% whole-market bootstrap interval of
[-0.08, 0.50] (Supplementary Table~\ref{tab:revision-band}).
This comparison changed width adaptation while retaining the response
to departures from a moving reference price. It therefore tested one
part of the price feedback rather than removing technical trading.

Additional comparisons used dealer depths $h=40,100,400$ at fraction $f=0.20$ with
active technical traders (Supplementary Table~\ref{tab:revision-impact}).
For the constant-product dealer, differentiating log price with respect to
normalized net trader purchases at initialization gives $2N_{\rm tr}/H$,
approximately $2/h$. The chosen depths therefore spanned initial slopes
0.05, 0.02, and 0.005. No independent exponential impact coefficient
was varied: changing depth also changed dealer funding. These
comparisons replaced an impact-coefficient sweep belonging to the
superseded standing-quote model. They measured sensitivity within the
corrected mechanism and did not hold available reserves fixed while
changing price sensitivity. Neither the quoted price nor the length of a step was
calibrated to an observed market.

The horizon campaign used fifty starting markets at fractions
0.140, 0.145, through 0.250, extending to 0.300 for grid size $L=32$ to bracket
short-horizon responses. Held technical traders prevented financial
feedback throughout. Endpoint comparisons used the original 3,000 steps
and size-dependent observation periods $3000(L/64)^z$ for $z=1,2$, alongside longer
observations. Majority criteria of 0.4, 0.5, and 0.6 were extracted
from each history. Interpolation and bootstrap intervals were calculated
separately for every size, horizon, and criterion
(Supplementary Table~\ref{tab:revision-horizons}). The no-switch
extension followed ten starting markets per technical-trader condition for
100,000 additional steps. The observed outcomes were active technical traders: 0/10 majorities, 95\% Wilson interval [0.000, 0.278]; held technical traders: 0/10 majorities, 95\% Wilson interval [0.000, 0.278].
These extensions reused starting markets and did not constitute
independent replications. The observed frequencies of Sell majorities described
the stated interval without establishing behavior at indefinitely long
times.

We also characterized the shared pre-switch histories at dealer depth $h=100$ with
active technical traders. Across the fifty starting markets, mean
interior Sell fractions were 0.2494 over
steps 501--750 and 0.2504 over steps
751--1,000. In the latter interval, the median within-market log-return
standard deviation was 0.000323, and
median within-market first and ninety-ninth percentiles were
-0.000730 and
0.000729. The mean fraction of exactly
zero computed log returns was 0.007.
These summaries described the starting markets without establishing
stationarity or matching a calibrated empirical return distribution.
The no-switch controls separately measured post-switch baseline
financial variation across dealer depths
(Supplementary Table~\ref{tab:paired-switch}).

The refined finite-size analysis resampled entire starting markets across
fractions in 2,000 bootstrap replicates. The empirical estimate used the
first upward crossing and allowed local decreases in probability, while
the logistic fit imposed a smooth increasing response
(Supplementary Table~\ref{tab:revision-fits}). Each replicate reselected
the susceptibility maximum over the tested fractions. Consequently,
intervals included variation in which sampled fraction supplied that
maximum, but not interpolation bias, uncertainty from selecting the
refinement grid, or changes from observing for longer. The recorded
valid-bootstrap fraction identified samples in which the requested
crossings remained bracketed. No susceptibility exponent or
infinite-size limit was fitted.
